% TOP_PROBLEM: {"id": 3150, "title": "Does the chromatic symmetric function distinguish between trees?", "category_id": 3, "category": {"id": 3, "name": "graph_theory", "display_name": "Graph Theory", "description": "Problems involving graphs, networks, and their properties.", "slug": "graph-theory", "order_index": 3, "created_at": "2026-07-31T15:26:25.671Z"}, "status": "open", "problem_number": "OPG-36880", "set_id": 12}
% FIRST_POSTED: 2026-10-01
% ATTEMPT_STATUS: unresolved
% UnsolvedMath ID 3150; source catalogue code OPG-36880
% !TeX program = lualatex
% Research attempt by GPT-6 Astra Ultra; no claim of a new complete solution.
\documentclass[11pt,a4paper]{article}
\usepackage[margin=25mm]{geometry}
\usepackage{fontspec}
\setmainfont{lmroman10-regular.otf}[BoldFont=lmroman10-bold.otf,
 ItalicFont=lmroman10-italic.otf,BoldItalicFont=lmroman10-bolditalic.otf]
\usepackage{amsmath,amssymb,mathtools}
\usepackage{unicode-math}
\setmathfont{latinmodern-math.otf}
\AtBeginDocument{\let\setminus\smallsetminus}
\usepackage{xurl,hyperref}
\hypersetup{colorlinks=true,urlcolor=blue}
\setcounter{secnumdepth}{0}
\setlength{\parindent}{0pt}
\setlength{\parskip}{5pt}
\setlength{\emergencystretch}{3em}
\begin{document}
\section*{Does the chromatic symmetric function distinguish between trees?}\label{3150}
{\small UnsolvedMath ID 3150 · OPG-36880 · Graph Theory · OpenGarden}
\subsection{Definitions and mathematical statement}
Let $G=(V,E)$ be a finite simple graph and let
$x_1,x_2,\ldots$ be commuting indeterminates. A proper colouring by
positive integers is a function $f:V\to\mathbb Z_{>0}$ with
$f(u)\neq f(v)$ whenever $uv\in E$. Define the chromatic symmetric
function by the formal sum
\[
 X_G(x_1,x_2,\ldots)
   =\sum_{f\text{ proper}}\prod_{v\in V}x_{f(v)}.
\]
This is a homogeneous symmetric function of degree $|V|$. Every
particular monomial has a finite coefficient, so no convergence
of the infinite sum is required. Its coefficient records the
number of proper colourings with the corresponding sizes of
colour classes.

A tree is a finite connected graph without a cycle. The question
asks whether two nonisomorphic trees can have the same chromatic
symmetric function. In its affirmative-distinguishability form,
the target assertion is
\[
 X_T=X_{T'}\quad\Longrightarrow\quad T\cong T'
 \qquad\text{for all finite trees }T,T'.
\]
Equality means coefficientwise equality in the same indeterminates;
isomorphism means a vertex bijection preserving adjacency. Equality
automatically forces the trees to have the same order, by
homogeneity. The question therefore concerns information beyond
the number of vertices. Specializing the first $q$ variables to
one and all others to zero counts proper $q$-colourings, but the
full symmetric function retains the separate colour-class sizes
discarded by that specialization.
\subsection{Short English statement}
Does the full chromatic symmetric function distinguish every pair of nonisomorphic finite trees?
\subsection{Research attempt}
\paragraph{Scope and process.}
This is a GPT-6 Astra Ultra research attempt dated 1 October 2026, using
primary-source browsing and elementary mathematical checks. The canonical
definition above is retained. Results below are restricted results or
reductions, not a claim of a new solution to the entire catalogue question.
No formal proof assistant or external mathematical referee checked this text.


\paragraph{Literature and plan.}
Martin, Morin and Wagner [R1] develop the power-sum expansion and extract
several tree invariants from it, including substantially stronger results
than the elementary families treated here. We use that expansion with
a complete proof of the relevant coefficient interpretation, then test
a tempting simplification: retaining only the component sizes arising
from deletion of a single edge. This distinguishes paths, stars and
double-stars, but an explicit pair shows why it cannot reconstruct
all trees. No claim of novelty is made for these restricted results.

\paragraph{Recovering the edge-cut data.}
Write $p_j=\sum_i x_i^j$ and $p_\lambda=\prod_jp_{\lambda_j}$ for a
partition $\lambda$. For $F\subseteq E(G)$ let $\lambda(F)$ be the
partition of $|V|$ formed by the component sizes of $(V,F)$, counting
isolated vertices. Inclusion-exclusion over edges whose endpoints
receive the same colour gives
\[
 X_G=\sum_{F\subseteq E(G)}(-1)^{|F|}p_{\lambda(F)}.
\]
Indeed, requiring every edge of $F$ to be monochromatic forces exactly
each connected component of $(V,F)$ to be monochromatic. Its freely
chosen colour contributes $p_j$ when that component has $j$ vertices.
Inclusion-exclusion removes every colouring that violates properness.
The argument is coefficientwise, so it needs no convergence assumption.

If $G=T$ is an $n$-vertex tree, every spanning subgraph is a forest,
and therefore $|F|=n-\ell(\lambda(F))$. All contributions to a fixed
power-sum coefficient have the same sign:
\[
 [p_\lambda]X_T=(-1)^{n-\ell(\lambda)}
 \#\{F\subseteq E(T):\lambda(F)=\lambda\}.
\]
Since the power sums are a basis over $\mathbb Q$, equality of symmetric
functions gives equality of all these integers. Partitions with exactly
two parts correspond precisely to deleting one edge of $T$.
Thus $X_T$ determines the multiset of unordered component-size pairs
produced by every single-edge deletion.

\paragraph{Three families that the data reconstruct.}
For $n\ge3$, the number of cuts of type $(n-1,1)$ is the number of
leaves. Deleting a leaf edge gives such a cut, and a singleton component
can only arise this way. At $n=2$ the sole cut has two singleton sides,
so the two leaves should instead be handled directly.

A tree of order $n\ge3$ with $n-1$ leaves is a star: its sole remaining
vertex must be adjacent to every leaf. A tree with precisely two leaves
is a path, since the degree sum gives
\[
 \#\{v:d(v)=1\}=2+\sum_{d(v)\ge3}(d(v)-2),
\]
so no vertex can branch. Hence $X_T$ determines these trees among all
trees, not just among previously specified stars or paths.

For $n\ge4$, a tree with $n-2$ leaves has exactly two internal vertices.
They are adjacent, because the path between them cannot have a leaf as
an intermediate vertex. Let $a,b\ge1$ be their respective numbers of
pendant neighbours. All edges except the central edge produce singleton
cuts; the central edge produces $(a+1,b+1)$. The single-edge cut multiset
therefore recovers the unordered pair $\{a,b\}$ and hence the double-star.
This also includes the four-vertex path without ambiguity.

\paragraph{An explicit obstruction to reconstruction from these cuts alone.}
On vertices $0,\ldots,6$, consider the two trees with edge sets
\[
 \begin{split}
 E(T_1)&=\{01,02,03,04,15,56\},\\
 E(T_2)&=\{01,02,03,14,15,26\}.
 \end{split}
\]
Their degree multisets are respectively $(4,2,2,1,1,1,1)$ and
$(3,3,2,1,1,1,1)$, so they are nonisomorphic. In both, the smaller
component sizes obtained from deleting the six edges are
$(1,1,1,1,2,3)$. Thus all two-part coefficients agree.
The full symmetric functions do distinguish them: the coefficient of
$p_{(3,2,2)}$ is zero for $T_1$ and one for $T_2$. For $T_2$, keeping exactly $03,14,15,26$ gives components
$\{0,3\}$, $\{1,4,5\}$ and $\{2,6\}$. This is the unique such
subset: leaves $4,5$ force the three-vertex component through $1$,
then leaves $3,6$ force the two remaining pairs. In $T_1$, leaves
$2,3,4$ would all need to be joined to $0$ to avoid singleton
components, producing a component of size at least four. Hence its
coefficient of this type is zero.

\paragraph{Verification and final status.}
The reproducible script
\texttt{scripts/check\_attempt\_batch02\_algebra\_2026\_10\_01.py}
enumerates all $2^6$ subsets of each displayed edge set, computes their
connected components including isolated vertices, and checks both the
cut-profile collision and the stated power-sum coefficients exactly.
The collision concerns a discarded part of the invariant; it is not
a counterexample to the catalogue question.
\textbf{UNRESOLVED.} The first remaining gap is to reconstruct a general
tree from the full multiset of component-size partitions. The two-part
projection has been shown insufficient by an actual pair, so a general
proof must use interactions among multiple deleted edges.

\subsection{Sources}
\begin{itemize}
\item[S1] ULAM.AI and the UnsolvedMath Contributors, supplied problems.json, record 3150 (OPG-36880), OpenGarden collection. Catalogue identity and source transcription; CC BY 4.0.
\url{https://huggingface.co/datasets/ulamai/UnsolvedMath}.
\item[S2] Open Problem Garden, Does the chromatic symmetric function distinguish between trees?. Original problem and discussion; the mathematical formulation below is expanded from the supplied source record.
\url{http://www.openproblemgarden.org/op/does_the_symmetric_chromatic_function_distinguish_trees}.

\item[R1] Jeremy L. Martin, Matthew Morin and Jennifer D. Wagner,
\emph{On distinguishing trees by their chromatic symmetric functions}.
Primary author-hosted paper, including the power-sum expansion and
reconstruction results; checked 1 October 2026.
\url{https://jeremymartinmath.github.io/papers/CSFTree.pdf}.

\end{itemize}
\end{document}
